\documentclass[11pt,a4paper]{article}
\usepackage[T1]{fontenc}
\usepackage{lmodern,amsmath,amssymb,amsthm,mathtools,booktabs,microtype}
\usepackage[margin=27mm]{geometry}
\usepackage[colorlinks=true,linkcolor=blue,citecolor=blue,urlcolor=blue]{hyperref}
\setlength{\emergencystretch}{2em}
\newtheorem{theorem}{Theorem}[section]
\newtheorem{lemma}[theorem]{Lemma}
\newtheorem{proposition}[theorem]{Proposition}
\newtheorem{corollary}[theorem]{Corollary}
\newcommand{\Q}{\mathbb Q}\newcommand{\Z}{\mathbb Z}
\DeclareMathOperator{\Nm}{N}\DeclareMathOperator{\Tr}{Tr}
\title{A quartic integral label and scalar norms from the detector field}
\author{}\date{16 November 2026}
\begin{document}\maketitle
\begin{abstract}
We construct a degree-four polynomial reflecting rational integrality
and a prepared scalar which is an explicit norm from
$K=\Q(\sqrt2)$, the field of the elliptic detector. The rational
argument of the scalar contact-height theorem is a fixed rescaling of
the integral elliptic label. This distinction retains the five-contact
height estimate while allowing the norm field to coincide with $K$.
The scalar has a uniformly positive odd valuation at the consumed prime,
square reduced denominator on integer inputs, and the required local
preparation at $2,3,17$. Under the retained arithmetic premises this
gives a twenty-five-variable defining family. Both prepared physical
weights can be chosen as norms from $K$ without additional prime-pattern
congruences. No realization with prescribed detector traces or smaller
variable count is asserted.
\end{abstract}

\section{Constants and statement}
We use the unsquared general source of~\cite{UnitNorm}, with the
elliptic subgroup, index-height, contact-height, pole-realization,
prime-pattern and aggregate trace-height statements in their retained
scopes~\cite{General,Dependencies}. Let $S$ be its exceptional set,
consisting of all primes at most $76176$ together with
$1875611,3828397,3982747$.

Put
\begin{equation}\label{eq:constants}
\begin{gathered}
 K=\Q(\sqrt2),\qquad n=116576341082536192,\qquad D=5951086353,\\
 \ell=76249,\qquad\eta=102,\qquad
 A=1296233=17\ell,\qquad Z_0=12556178240.
\end{gathered}
\end{equation}
The integer elliptic label and the rational scalar contact argument are
\begin{equation}\label{eq:label}
 Z(X)=AX^2+Z_0,\qquad
 \beta(X)=2Z(X)^2-D^2,\qquad q(X)=\beta(X)/n.
\end{equation}
In particular
\[
 \beta(X)=3360439980578X^4+65102930354279680X^2
                              +279899795208436154591.
\]
Define
\begin{equation}\label{eq:contacts}
 R(X)=\beta(\beta-n)(\beta-2401n)(32\beta+343n)
        =n^4r_4(q),\qquad r_4(t)=t(t-1)(t-2401)(32t+343).
\end{equation}
Let $M=R(0)$, namely
\[
\begin{split}
M={}&189711668070091549022366223774716198062888712249407876629023195358725017888.
\end{split}
\]
For every fixed odd integer $L\ge3$, put
\begin{equation}\label{eq:scalar}
 C_L=\eta^2M(\ell^{(L+1)/2}+1),\qquad
 \rho_L(X)=\frac{\eta^2\ell^LMR(X)}{(1+C_LR(X))^2}.
\end{equation}

\begin{theorem}\label{thm:main}
Under the retained arithmetic premises, use the label $\beta$ in the
comparison tests and the scalar $\rho_L$ in the unsquared general
source. This gives a computable family of integer polynomials in
twenty-five total variables. Every sufficiently large fixed odd $L$
defines $\Z$ over $\Q$. Completeness holds for every integer input
and every odd $L\ge3$.

For every rational input, $\rho_L$ is a nonzero norm from $K$.
For each integer input, the prepared witnesses can be chosen so that
both physical weights $w_0,w_1$ are norms from $K$. Consequently both
conics
\[
                         u_j^2-2y_j^2=4w_j/17\qquad(j=0,1)
\]
have rational points. These assertions do not prescribe their first
coordinates. No numerical sufficient value of $L$ or degree bound for
the resulting defining determinant is asserted.
\end{theorem}

\section{Norm identities and integrality}
Use the exact four-square certificate
\begin{equation}\label{eq:squares}
\begin{gathered}
 (\theta_0,\theta_1,\theta_2,\theta_3)=(0,1,2401,-343/32),\\
 (a_0,a_1,a_2,a_3)
 =(5951086353,5960872849,17757117551,5845158349),\\
                         D^2+n\theta_i=a_i^2\quad(0\le i\le3),
\end{gathered}
\end{equation}
also used in~\cite{NormPreparation}. Define
\[
 \Gamma(X)=4(2+\sqrt2)\prod_{i=0}^3(a_i+Z(X)\sqrt2).
\]
Since $\Nm(4(2+\sqrt2))=32$ and
$\Nm(a_i+Z\sqrt2)=-(\beta-n\theta_i)$, multiplication gives
\begin{equation}\label{eq:normR}
                              \Nm\Gamma(X)=R(X).
\end{equation}
The identities
\[
 \Nm(357+160\sqrt2)=\ell,\qquad \Nm(5+2\sqrt2)=17
\]
then give the explicit scalar norm
\begin{equation}\label{eq:normrho}
 \rho_L(X)=\Nm\left(
   \frac{\eta(357+160\sqrt2)^L\Gamma(0)\Gamma(X)}
        {1+C_LR(X)}\right).
\end{equation}

\begin{lemma}\label{lem:integrality}
The polynomial $\beta$ has integer coefficients and degree four, and
\[
                  \beta(X)\in\Z\quad\Longleftrightarrow\quad X\in\Z
                  \qquad(X\in\Q).
\]
For every real $X$ one has $\beta(X)\ge\beta(0)>2401n$ and
$R(X)\ge M>0$.
\end{lemma}
\begin{proof}
Write $\beta=A_4X^4+A_2X^2+A_0$, where
$A_4=2A^2$, $A_2=4AZ_0$, and $A_0=2Z_0^2-D^2$.
The leading coefficient has valuation $1$ at $2$, $2$ at
$17,\ell$, and $0$ at every other prime. The quadratic coefficient
has valuation $8$ at $2$, $1$ at $\ell$, and $2$ at $17$;
it is integral elsewhere.

If $v_p(X)=-e<0$, the leading term therefore has valuation
$1-4e$ at $2$, $2-4e$ at $17,\ell$, and $-4e$ elsewhere.
These are negative and strictly smaller than the respective lower
bounds $8-2e$, $1-2e$ at $\ell$, $2-2e$ at $17$, and $-2e$
elsewhere for the quadratic term. They are also smaller than the
valuation of the integral constant. Thus the leading term uniquely
dominates and $\beta(X)$ is nonintegral. This includes $p=2,e=1$,
when the leading valuation is $-3$. The converse is immediate.

For the real statement, $Z(X)\ge Z_0>0$ and
$\beta(0)-2401n=269266757599>0$. All four factors of $R$ are
positive and increasing on this label range.
\end{proof}

\section{Scalar preparation}
\begin{lemma}\label{lem:preparation}
For every real input,
\[
               0<\rho_L(X)<\frac1{\eta^2\ell MR(X)}<1.
\]
At the consumed prime,
\[
 v_\ell(\rho_L(X))=
 \begin{cases}
 L,&v_\ell(X)\ge0,\\
 L-8+16e,&v_\ell(X)=-e<0.
 \end{cases}
\]
These valuations are positive and odd. For integer inputs, $\rho_L$
has square reduced denominator, is a square in
$\Q_2,\Q_3,\Q_{17}$, and has exact valuations $(12,2,2)$ there.
\end{lemma}
\begin{proof}
The real bound follows from $R\ge M>0$ and
$C_L^2>\eta^4M^2\ell^{L+1}$. In particular the denominator in
\eqref{eq:scalar} and~\eqref{eq:normrho} never vanishes.

The exact residues modulo $\ell$ are
\[
 Z_0\equiv26663,\qquad \beta(0)\equiv14480,\qquad
 M\equiv68886,\qquad 1+\eta^2M^2\equiv46270.
\]
For an $\ell$-integral input, $\ell\mid A$ gives $R(X)\equiv M$,
and $C_L\equiv\eta^2M$; both $R$ and the denominator are units.
For $v_\ell(X)=-e<0$, Lemma~\ref{lem:integrality} gives
$v_\ell(\beta)=2-4e$ and $v_\ell(R)=8-16e<0$.
As $C_L$ is an $\ell$-unit, the denominator has this latter
valuation. The stated formula follows.

Now let $X\in\Z$. The label is odd.
Since $n=256N$ with $N$ odd,
\[
 R/32=\beta(\beta-n)(\beta-2401n)(\beta+343n/32)
\]
is odd and congruent to $\beta^4\equiv1$ modulo eight. Hence
$v_2(R)=v_2(M)=5$ and $R/M\in1+8\Z_2$.
At three, $A\equiv Z_0\equiv2$, so $Z$ is congruent to two or
one according as $3\mid X$ or $3\nmid X$. Since $n\equiv1$
and $D\equiv0$, in both cases
$\beta\equiv2$ and all four factors of $R$ are units.
Thus $R/M\in1+3\Z_3$.
At seventeen, $A,Z_0$ are divisible by $17$,
$\beta\equiv13$, $n\equiv15$, and $R\equiv M\equiv13$.
The nonconstant label coefficients are divisible by $17^2$, so
$R/M\in1+17\Z_{17}$.

These three principal-unit groups consist of squares. The prime
$\ell$ is also a local square, since its residues modulo
$8,3,17$ are $1,1,4$. The denominator is a unit at all three
places because $2\cdot3\cdot17\mid C_L$. Writing
$\rho_L=\eta^2\ell^LM^2(R/M)/(1+C_LR)^2$ proves the square
assertions and the exact valuations $(12,2,2)$.

Finally the displayed scalar numerator and denominator are relatively
prime on integer inputs. Any prime dividing $\eta M$ divides $C_L$;
a prime dividing $R$ cannot divide $1+C_LR$; and the residue
certificate excludes $\ell$ from the denominator. The reduced
denominator is consequently a square.
\end{proof}

\section{Transfer through the general source}
Retain all nine point pairs, sixteen comparisons, six scalar tests,
local orders and contact sets of~\cite{UnitNorm}. The target label
in its comparison $L_1-\beta$ is~\eqref{eq:label}. For clarity, use
$\mathcal R_j$ for the source contact functions, so
\[
 a_0=1,\quad a_1=1+\rho_LT^2,\quad t_j=a_j\lambda,
 \quad\delta_j=t_j-\gamma_j,\quad
 \mathcal R_j=\frac{\prod_{\theta\in\Theta_j}(t_j-\theta)}
                         {\delta_j^{10}},\quad w_j=\rho_L\mathcal R_j.
\]
The detector products $\mathcal W_j$ and inputs
\[
                         b_j=4w_j/\Tr(\mathcal W_j^{-1})^2
\]
are unchanged apart from the stated label and scalar substitutions.

\begin{proof}[Proof of the defining-family assertion]
Let $S'$ be $S$ together with the prime divisors of $M$.
It is fixed and excludes $\ell$. The factorization
\[
                  n=2^8\cdot11\cdot13\cdot101\cdot757\cdot5939\cdot7013
\]
puts every prime divisor of $n$ in $S$. Outside $S'$, any finite
odd contact of the rational argument $q=\beta/n$ gives a positive
odd valuation of $\rho_L$: precisely one factor of $R$ has positive
odd valuation, all fixed multipliers are units, and the normalizing
denominator is a unit. Negative valuations of $q$ outside $S$ are
even by Lemma~\ref{lem:integrality}; thus the infinity contact has
no odd contribution there. The nonzero residue of $R$ at $\ell$
excludes a finite contact at the consumed prime.

The retained five-contact height theorem is stated for every rational
argument, so it applies to $q$, not only to the elliptic label.
Its scalar-odd consequence, with the unsquared-scale proof
of~\cite{UnitNorm}, gives
\[
 h(q)\le Hc^{7/9},\qquad
 h(\beta)\le h(q)+\log n\le H'c^{7/9}.
\]
Here $c=|a_{\rm root}-1|>0$ is the absolute adjacent root-index difference of the source. Omitting the fixed
primes dividing $M$ changes only the constants. At the consumed prime,
Lemma~\ref{lem:preparation} and the same scalar-odd lemma give
\[
                        v_\ell(c)\ge2L+1,\qquad c\ge\ell^{2L+1}.
\]
Since $R$ is a degree-four polynomial in the label, the elementary
height estimate for~\eqref{eq:scalar} is
\[
                 h(\rho_L)\le12h(\beta)+2L\log\ell+O(1),
\]
where the implied constant is fixed. All local comparison extraction,
integer residual divisibilities, aggregate trace-height bounds and
Lucas growth inequalities are otherwise unchanged. Their strict
exponents consequently force $\beta(X)\in\Z$ for sufficiently large
fixed odd $L$, and Lemma~\ref{lem:integrality} gives $X\in\Z$.

For completeness, an integer input gives a positive odd integer label
and a fixed positive prepared scalar. Follow the source's order of
choices. First fix the scalar point and the adjacent root index so
that all six scalar functions are sufficiently small at the finite
support of $\rho_L$, at the fixed exceptional primes, and in the
real closure. Then set the first root label to $\beta(X)$ and choose
the other three labels by the same Lucas construction. All labels
are positive and odd. Fix the permitted auxiliary congruence class,
real neighborhoods, denominator avoidance and bounds for comparisons.

The rational-contact prime-pattern theorem applies to the same fixed
coefficients $-\theta/a_j$. Include the finite prime divisors of $M$
in the avoidance conditions before prime selection. The theorem gives
the required denominator divisibility, arbitrarily large positive
slope, Gaussian reduction orders and common auxiliary-index refinement.
Both $\mathcal R_j$ tend to one as the slope tends to infinity at
each preparation place, so the same scalar and comparison-factor
preparation applies. The extra divisibility of the auxiliary index by
the selected primes still follows from their prime-cluster bound.
It preserves the scalar normalization denominators and every
comparison precision in exactly the source's calculation.

The reduced denominators of the two pole inputs are square by the
unchanged prime-by-prime argument. At $2,3,17$ their squareclasses are
those of $\rho_L$, and their valuations are $(12,2,2)$ once the
products and contact functions are sufficiently close to one.
These meet the source pole-preparation bounds. At the real places
$1/2<\mathcal R_j<2$, $1\le\mathcal W_j<2$ and
$0<b_j<8\rho_L<17$. Thus the inherited pole-realization theorem
supplies both rational pole witnesses.

All point and difference guards follow from nonzero labels and the
same refined auxiliary-index conditions. The scalar denominators and
the whole finite-algebra guards follow from the uniform neighborhoods,
positivity, and branch calculations of~\cite{UnitNorm}; these use
the stated scalar preparation and positive labels, rather than an
identity between the label and its scalar contact argument. The
retained-coordinate census is unchanged. The finite-algebra
construction therefore produces a computable defining family in
twenty-five variables.
\end{proof}

\section{Prepared norms in the detector field}
\begin{lemma}\label{lem:prime}
Every rational prime $p\equiv1\pmod8$ is a norm from $K$.
\end{lemma}
\begin{proof}
At $p$, quadratic reciprocity gives $(2/p)=1$, so the extension
is split. At two, $p$ is a local square and hence a norm. At every
other finite place the extension is unramified or split and $p$ is
a local unit, hence a norm. The real extension is split. The conic
Hasse principle~\cite[Theorem~26.1 and Corollary~26.5]{Poonen}
therefore gives a rational solution of $x^2-2y^2=pz^2$.
Its coordinate $z$ is nonzero because $2$ is not a rational square.
\end{proof}

\begin{proof}[Proof of the norm assertions in Theorem~\ref{thm:main}]
Equation~\eqref{eq:normrho} proves the scalar assertion. The retained
Gaussian prime pattern already gives every contact prime congruent to
one modulo eight~\cite[Section~3]{Dependencies}. For the prepared
slope $\lambda=u/v$, the exact contact identity is
\[
 \mathcal R_j=\left(\frac{a_j}{v\delta_j}\right)^{10}
               \prod_{\theta\in\Theta_j}
                           \bigl(u-(\theta/a_j)v\bigr).
\]
Every product factor is a selected prime and hence a norm by
Lemma~\ref{lem:prime}; the prefactor is a rational square. Thus both
$\mathcal R_j$ and $w_j$ are norms. As $17=\Nm(5+2\sqrt2)$
and $4=\Nm(2)$, both bare conics in the theorem have rational points.
\end{proof}

\begin{corollary}\label{cor:conditional}
In either unchanged thirteen-variable conditional source
of~\cite{Conditional,Direct}, prepared
physical weights can be chosen as norms from $K$, with no additional
prime-pattern congruences. Both conics $u_j^2-2y_j^2=4w_j/17$ have
rational points. No prescribed-trace realization follows.
\end{corollary}
\begin{proof}
The conditional source has six contacts and physical weights
$w_j=F_j/(4\delta_j^8)$, with $t_j=a_j\lambda$ and
$\delta_j=t_j-1$. For its prepared slope $\lambda=u/v$,
\[
 w_j=\left(\frac{a_j^3}{2v^3\delta_j^4}\right)^2
             \prod_{\theta\in\Theta}\bigl(u-(\theta/a_j)v\bigr).
\]
The factors are the same permitted positive prime values, already
one modulo eight. Apply Lemma~\ref{lem:prime} and the certified norm
of seventeen.
\end{proof}

For a reduction to twenty-one or nine variables, one must still
realize the actual simultaneous equations
\[
       \bigl(m_j\Tr(\mathcal W_j^{-1})\bigr)^2-d_jy_j^2=4w_j/17
\]
with admissible levels and square classes and with all inherited
guards. The norm identities above do not prescribe these first
coordinates and do not supply that additional arithmetic theorem.

\section{Exact verification}
The accompanying standard-library program
\begin{center}\small
\nolinkurl{computation/pole_parameters/detector_field_norm_v52/check.py}
\end{center} verifies the full four-square
certificate, reconstructs the quartic label, expands the norm identity
as an exact polynomial identity, checks all fixed norms and residue
certificates, and verifies the coefficient divisibilities used at
$2,3,17$. Its certificate includes every coefficient of the
degree-sixteen contact polynomial. Explicit rational denominator
examples supplement, but do not replace, the uniform valuation proof.
The source-transfer proof uses the retained arithmetic theorems and is
not a conclusion of finite experimentation.

\begin{thebibliography}{9}\raggedright
\bibitem{UnitNorm}\emph{A single norm determinant for the 25- and
13-variable constructions}, project manuscript, 14 November 2026,
\nolinkurl{papers/finite_algebra/unit_norm_20261114_v28.tex}.
\bibitem{General}\emph{A defining polynomial in twenty-five variables:
combined scalar and comparison conditions}, project manuscript,
10 November 2026,
\nolinkurl{papers/few_variables25/few_variables25_merged_20261110_v19.tex}.
\bibitem{Dependencies}\emph{Arithmetic dependencies of the construction
with 29 variables}, project manuscript, 8 November 2026,
\nolinkurl{papers/arithmetic_dependencies/arithmetic_dependencies_20261108_v8.tex}.
\bibitem{NormPreparation}\emph{Norm preparation for general and
conditional elliptic detectors}, project manuscript, 16 November 2026,
\nolinkurl{papers/pole_parameters/norm_compatible_scalar_20261116_v47.tex}.
\bibitem{Conditional}\emph{Thirteen variables from a weak $abc$
inequality}, project manuscript, 9 November 2026,
\nolinkurl{papers/conditional_few_variables/conditional13_20261109_v16.tex}.
\bibitem{Direct}\emph{Thirteen variables from direct elliptic index
comparison}, project manuscript, 10 November 2026,
\nolinkurl{papers/conditional_few_variables/direct_input13_20261110_v22.tex}.
\bibitem{Poonen} B.~Poonen, \emph{Introduction to Arithmetic Geometry},
MIT 18.782 notes, 2009, Theorem~26.1 and Corollary~26.5,
\url{https://math.mit.edu/~poonen/782/782notes.pdf}.
\end{thebibliography}
\end{document}
